\cotefile{83}
\scanpage{96}
Supposons maintenant que
\[
\delta=\delta'=0\qquad \text{i.e.}\qquad b^2-4c=b'^2-4c'=0
\]
\unsure{i.e.} ce qui signifie \deleted{que \unsure{si on considère} $2$ alg. d'Az.} \inserted{que $A_1,A_2$ sont} loc. isom. à $k[T]/T^2$

% slanted
\margin{\uuline{faux} si on ne suppose pas que \unsure{l'on} est sur un corps, \uuline{ou} \unsure{$2$ inv.}}

(je l'ai peut-être vérifié dans le temps -- ça n'a pas l'air \deleted{évident} : fait évident, sur une base générale -- évident sur un corps, plus gén. si $\bigl[$On est ramené à ceci : si $b\in k$ tel que $b^2=0$, $\exists$ loc. fppf $\alpha,\beta\in k$, $\beta\in k^*$, tels que $\alpha^2=0$, $\beta(2\alpha+b)=0$ Evident si $2$ inv., mais faux si $2.1_k=0$ (car. $2$) car \deleted{\unsure{$\beta b=$|$\beta b=0$}} alors l'équation en $\beta$ donne $\beta b=0$ donc (comme $\beta\in k^{\unsure{*}}$) $b=0$, mais des situations $b^2=0$, $b\ne0$ ça \unsure{court|courre|croule|coince} ! $\bigr]$

\deleted{Ça signifie ainsi que} \deleted{Mais des alg. d'Az} \deleted{$A_1,A_2$ sont des sous-algèbres \unsure{correspondant|conjuguées}} \deleted{\inserted{normalis.}} \deleted{à la partie unipotente} \deleted{tirées d'un $\Theta\in M$ (Azumaya ambiant) qui vér.} \deleted{de Trace}
\[
\deleted{\operatorname{Tr}\Theta=\det\Theta=0}
\]
\deleted{Les \unsure{voici les normalisateurs|les unipotents} d'un tel} Il y a corr. $1-1$ entre une
% archive sheet number: 95
\scanpage{97}
Donc on s'intéressera \inserted{finalement} au cas où $A_1=k.1+D_1$, $A_2=k.1+D_2$, $D_1,D_2$ contenus dans $M^0=\{x\in M\mid\operatorname{Tr}x=0\}$ et isotropes i.e.
\[
\begin{aligned}
\operatorname{Tr}u&=\operatorname{Tr}v=0\\
\det u&=\det v=0
\end{aligned}
\]

Dans ce cas là on trouve
\[
Q(\unsure{t|b},b,b',c,c')=t^2
\]
Mais dans ce cas là on vérifie aussitôt que pour $x\in M$
\[
\left[\begin{aligned}
\operatorname{Tr}(u\deleted{\unsure{v}}x)=0&\Longleftrightarrow x\text{ normalise }\deleted{\unsure{D|k}}ku\\
\text{i.e.}\qquad [x,u]&=xu-ux=0
\end{aligned}\right.
\]
\deleted{donc} \inserted{i.e.} $x\in B_u$\quad sous-\uuline{alg.} de Borel normalisatrice de $u$,

donc si $t=\operatorname{Tr}(uv)=0$ cela implique que $u,v$ engendrent \deleted{\unsure{un}} une sous-alg. de Borel
\[
k\oplus ku\oplus kv=B_u=B_v
\]
(et \unsure{inversement|nécessairemt}).

Donc
% archive sheet number: 96
\scanpage{98}
\uline{Proposition} Si $b=c=b'=c'=0$, alors $M[b,\deleted{\unsure{b}}c,b',c';t]=M(0,0,0,0,t)$ est Azumaya, sss \deleted{$t\in k^*$.} $Q(0,0,0,0,t)=t^2\in\deleted{\unsure{b}}k^\times$ i.e. $t\in k^\times$

En effet, si $t\in k^*$, alors on ``réalise'' $t$ dans l'algèbre de matrices $M_2$ p. ex. par $u=\begin{pmatrix}0&1\\0&0\end{pmatrix}$ et $v^{\unsure{\alpha}}=t\begin{pmatrix}0&0\\1&0\end{pmatrix}$ on trouve $\operatorname{Tr}uv=t$ OK.

Si $t$ \unsure{non} inv., \unsure{la|le} \deleted{\unsure{qui}} précède montre que ça ne se réalise pas dans Azumaya.

% slanted
\margin{Il faudrait ici montrer \inserted{(\unsure{déduire du cas particulier|éclaircir ce cas particulier})} que ce critère, en termes de $Q(\text{---})\in k^\times$, est valable sans supposer $u^2=v^2=0$ (avec la seule condition \deleted{\unsure{si} $\delta\delta'\in k^{\unsure{?}}$} \deleted{$\delta=0,\delta'=0$} $\delta\delta'=0$ -- de ttes façons cela se réduit au cas d'un corps par \unsure{changt} de base}

Reste le cas où $\deleted{\unsure{b}}\delta^{\unsure{2|'}}\ne0\quad\delta'=0$

\deleted{\unsure{Ici}} On peut supposer [un peu court !!]
\[
b'=c'=0
\]
donc voici encore
\[
Q(b,c,b',c',t)=t^2
\]
donc $Q$ inv. signifie $t$ inv.

Montrons que cela signifie que $M(b,c,b',c';t)$ est Azumaya.
% archive sheet number: 97
\scanpage{99}
\uline{Th} Soient $b,c,b',c',t\in k$, d'où $M=M(b,c,b',c',t)$, avec base can. $\underline{1},\underline{u},\underline{v},\underline{u}\,\underline{v}$, et formes $\operatorname{Tr}$ et $\det$ \ldots{} Pour que $M$ soit Azumaya, il faut et suffit que la forme $\det$ sur $\unsure{1k|k}\oplus k\underline{u}\oplus k\underline{v}$ soit non dégénérée, i.e. que son discriminant \unsure{arithmét|réduit} \unsure{et|ci} (précédemment désigné par $-Q$
\[
\Delta(b,c,b',c',t)=-(t^2-bb't+b^2c'+b'^2c-4cc')
\]
soit inv.

En termes plus intrinsèques, soit
\[
\deleted{M\ni}\qquad M\supset A_1,\deleted{\unsure{B}} A_2
\]
une algèbre \inserted{ass. unit. loc.} libre de rang $4$ sur $k$, \deleted{\unsure{sous une algèbre}} et deux sous-alg. $A_1,A_2$ \inserted{loc. libres} de rang $2$ \inserted{transverses}. On suppose que \deleted{\unsure{l'on peu}} dans $M$ il y ait des formes $\operatorname{Tr},\det$ avec les propr. hab, \inserted{et que $A_1,A_2$ ne \uuline{sont} pas (sur aucune fibre) contenues dans une sous-alg. de rg $3$ ($A_1+A_2$).} (ce qui implique que si on \unsure{décrit} sur $A_1$ gén. $u$, $A_2$ gén. $\unsure{v|V_2}$, d'où $b,\unsure{c|k},b',c',t$, on a $M\simeq M(b,c,b',c';t)$). Alors $M$ est azumaya sss la restriction de la forme \unsure{quad.|quadr.} $\det$ sur $A_1+A_2$ est non dégénérée,
% archive sheet number: 98
\scanpage{100}
\uline{Dém} C'est géom., on peut se placer sur corps alg. clos. Nécessité : Supposons que ce soit Azumaya, \deleted{On peut d'} \inserted{prouvons non dégénérescence.}

Si $A_1,A_{\unsure{1|2}}$ sont unipotentes, on a prouvé le résultat, et aussi si $A_1,A_2$ sont toutes les deux étales. Il reste le cas où l'une est étale, l'autre unipotente, mais il n'y a qu'une seule position relative \inserted{possible}. Dans ce cas on choisit $v$ nilpotent, $u$ un projecteur
\[
\begin{aligned}
&\deleted{b=1,\ c=0},\quad b'=0,\ c'=0\\
&Q=t^2,\quad\Delta=-Q=-t^2
\end{aligned}
\]
: vérifier que $t\ne\unsure{0|10}$. On choisit
\[
\deleted{\begin{pmatrix}0&1\\0&0\end{pmatrix}=v,\quad u=\begin{pmatrix}1&0\\0&0\end{pmatrix}\quad uv=\begin{pmatrix}0&1\\0&0\end{pmatrix}}
\]
\[
\begin{aligned}
v&=\begin{pmatrix}0&1\\0&0\end{pmatrix}&u&=\begin{pmatrix}\alpha&\beta\\\gamma&\delta\end{pmatrix}\quad\text{``général'',}\\
uv&=\begin{pmatrix}0&\alpha\\0&\gamma\end{pmatrix}&\operatorname{Tr}uv&=\gamma\\
&&\delta(\unsure{u|x})&=(\alpha+\delta)^2-4(\alpha\delta-\beta\gamma)
\end{aligned}
\]
On peut choisir $\alpha\beta\gamma\delta$ tels que $\gamma\delta(\alpha,\beta,\gamma,\delta)\ne0$

Suffisance \deleted{Il suffit de} \ldots.
% archive sheet number: 99
\scanpage{102}
\uline{Azumaya rang 4}\hfill (II)\hfill Avril 86

Soit $A$ Alg. d'Azumaya rang $4$ sur \deleted{schéma $S$} \inserted{anneau \inserted{$k$ (disons)}}, et $B\subset A$ un sous-\deleted{Module loc. \unsure{ct.|f.} direct} \inserted{fibré vectoriel} de rang $3$. Alors $B^\perp$ est un sous-\deleted{\unsure{vectoriel} de} \inserted{fibré vectoriel} de rang $1$, soit $D$. On a corr. $1$-$1$ entre les $B$, et les $D$, car \unsure{les formes quadratiques|la forme quadratique} canonique $\operatorname{Tr}(u,v)$ sur $A$ est non dégénérée.

À vrai dire, il y a deux formes bil. sym. non dégénérées intéressantes sur $A$, savoir
\[
\varphi(u,v)=\operatorname{Tr}(uv)\qquad\text{et}\qquad\inserted{\psi(u,v)=}\varphi_{\det}(u,v)=\operatorname{Tr}u\operatorname{Tr}v-\operatorname{Tr}(uv),
\]
\deleted{donnant} \deleted{ayant des discriminants \unsure{qui ne}}.

Je dis que pour tt fibré vectoriel $B$ (de rang quelconque) qui contient $1$, \uline{les deux formes ont} $=$ \uline{discriminant sur} $B$. En effet, il suffit de le voir dans cas universel d'une grassmannienne sur $\operatorname{Spec}\mathbb{Z}$, ce qui nous permet de supposer $2$ inversible, donc
\[
A=k.1\oplus A^{\unsure{s|\leftrightarrow|\natural}}
\]
\margin{$A^s\overset{\text{déf}}{=}\{u\in A\mid\operatorname{Tr}u=0\}= \mathbb{1}^{\perp}$ (\uline{NB} L'orthogonal de $\mathbb{1}$ pour l'une ou l'autre forme c'est le $=$ !)}
Donc
\[
B=k.1\oplus V,\quad V\subset A^s
\]
% archive sheet number: 101
\scanpage{103}
et \deleted{\unsure{l'|$\xi$}} on est ramené à prouver l'assertion \deleted{pour} séparément pour $k.1$, et $V$ (égalité des discr.). Or sur $k1$ les deux formes \deleted{bil.} \unsure{qua. ass.|qu. ass.} prennent la $=$ valeur $2$, OK. Sur $V$, \deleted{c'est} elles sont égales ! On gagne !

\deleted{\uline{Proposition}}

\uline{NB} Soit $B$ sous-module de $A$ tel que $B\supset1$, alors
\[
B^{\perp(\varphi)}=B^{\perp(\psi)}.
\]
En effet, \deleted{les deux orthogonaux} \deleted{c'est vrai $B^\perp$ est formé des} on a
\[
\varphi(1,u)=\psi(1,u)=\operatorname{Tr}u
\]
donc les orthogonaux de $1$ sont les $=\text{s}$ (c'est $A^{\unsure{s|\leftrightarrow|\natural}}=\{u\in A\mid\operatorname{Tr}u=0\}$), et $B^\perp\subset A^s$. Or si $\deleted{\unsure{x}}\zeta\in A^s$, on a pour \deleted{\unsure{ql|$\varphi$}} \deleted{tout} $\deleted{\unsure{\xi}}u\in A$
\[
\psi(u,\zeta)=-\varphi(u,\zeta)
\]
donc dire que $\zeta$ est orth. à $B$ pour $\varphi$, ou pour $\psi$, est pareil, cqfd.
% archive sheet number: 102
\scanpage{104}
On peut être plus explicite encore

\deleted{Plaçons} Les $\delta_B$ pour $B$ variable forment une section du fibré $\underline{O}_P(\unsure{2|1})\inserted{=\underline{O}_P(1)^{\otimes2}}$ sur le fibré projectif (de dim relative $2$)
\[
\mathbb{P}(A/k)
\]
des (hyper)plans dans $A$, donc proviennent d'un élémt canonique de
\[
\begin{aligned}
\Gamma(\underline{O}_P(2))&\simeq\operatorname{Sym}^2_k(A/k)=\\
&\simeq\text{espace des formes quadratiques}\\
&\hspace{1.4em}\text{sur}\quad(A/k)^\vee\simeq A^{\unsure{s|\leftrightarrow|\natural}}
\end{aligned}
\]

Or il y a une telle forme quadratique canonique sur $A^s$, savoir $\det$. Et on vérifie que $\delta_B$ est bien défini par cet élémt.

(Cela n'a rien à voir avec la structure multiplicative de $A$, qui n'intervient pas pour définir les $\delta_B$, il suffit de la forme quadratique \inserted{non dég.} $\det$ et l'élémt $1$ de $A$ \deleted{\unsure{non isotro|non injection}} \ldots
% archive sheet number: 103
\scanpage{105}
\[
\mu:\underbrace{\bigwedge^2 B/k}_{\substack{\parallel\\\det B/k}}\longrightarrow A/B
\]

Or on a un iso. can.
\margin{attention à conv. de signes}
\[
\det B\simeq\underbrace{\det k}_{k}\otimes\det B/k\simeq\det B/k
\]
donc on peut interpréter $\mu$ comme
\[
\mu:\underset{\substack{\rotatebox{90}{$\simeq$}\\(A/B)^{\otimes-1}}}{\det B}\longrightarrow A/B
\qquad\deleted{(\unsure{\simeq}\det B^{\unsure{-\otimes1|-\otimes}})}
\]
soit encore
\[
\mu\in A/B^{\otimes2}\simeq D^{-\otimes2}
\]

Ceci dit, il est clair que
\[
B\text{ sous-algèbre}\Longleftrightarrow\mu=0
\]
Il suffit de voir que
\[
\mu=\pm\delta
\]
\margin{le signe dépendant des conventions de signes}

Or il suffit de le voir dans le cas universel $S=\mathbb{P}^2_{\mathbb{Z}}$, et comme \unsure{nos|mes|ces} \deleted{il n'y a pas d'éléts} $\mu$ et $\delta$ \deleted{sont les $=\text{s}$} universels ont les $=\text{s}$ \deleted{zéros, \unsure{sans}} ens de zéros, on gagne !
% archive sheet number: 104
\scanpage{106}
laquelle disparait quand on passe au carré
\[
(\det B)^{\otimes2}\otimes\deleted{\unsure{A/B}}(A/B)^{\otimes2}\simeq1
\]
qui donne
\[
(\det B)^{-\otimes-2}\simeq(A/B)^{\otimes2}
\]
Ainsi, on peut écrire
\[
\delta_{\det\mid B}\in(A/B)^{\otimes2}\simeq D^{\otimes-2}
\]
[et en fait, c'est simplemt donné par la restriction de la forme quadratique $\det$ à $D$ : $\delta_{\det\mid B}=\det_{\mid D}$]

Or on a une application bilin can
\diagram{\begin{tikzcd}[ampersand replacement=\&,row sep=small]
\deleted{\unsure{\mu:}}\ \deleted{\unsure{A}}B\times B \arrow[r] \& A/B\\
(u,v) \arrow[r] \& uv\bmod B
\end{tikzcd}}
et celle ci est \unsure{alterné|alternée}, puisque par H.C, on a pour $u\in B$
\[
u^2\equiv0\qquad\bmod B
\]
(\deleted{\unsure{car}} $u^2=\operatorname{Tr}(u)u+\det u.1$, et $1,u\in B$).

Ainsi, on trouve
\[
\bigwedge^2 B\xrightarrow{\mu} A/B
\]
\unsure{de|Or|On} sorte que
\[
\mu(u\wedge v)=uv\bmod B.
\]
D'ailleurs \deleted{\unsure{$\mu$}} la fonction bil. $uv\bmod B$ est nulle quand l'un ou l'autre des \unsure{argument|arguments} est dans $k$, donc on peut considérer $\mu$ comme
% archive sheet number: 105
\scanpage{107}
correspondant
\[
\begin{aligned}
B_D=D^\perp\bigl(&=\text{normalisateur de }D\text{ dans}\\
&\text{l'alg. de Lie }A\text{ i.e.}\\
&\text{ens des }\unsure{u|x}\in A\text{ tels que}\\
&\deleted{[u,x]\in D\ \text{pour }\unsure{\forall|\text{tt}}\,x\in D}\\
&\operatorname{ad}_{\text{Lie}}(D)\subset D\bigr)
\end{aligned}
\]
\[
(B_D^\times=\text{normalisateur de }D\text{ dans }A^\times)
\]

D'où équivalence de (i) avec (iii$'$). D'autre part, (iii$'$) $\Rightarrow$ (iii) trivialemt, il reste à voir que (iii) $\Rightarrow$ (i) -- finalement je ne suis pas si sûr que c'est vrai, sauf sur un corps (ou un anneau réduit) le cas-test décisif est celui où $k=k_0[t]/t^{\unsure{2}}$ \ldots{} (histoire de cohomologie d'une algèbre de Borel \ldots)

\deleted{\uline{Corollaire}.} Notons que $\delta_B=\operatorname{disc}^{\prime}_{\det\mid B}$ est une section de
\[
\delta_B\in(\det B)^{\otimes-2}\left[\underset{\text{iso can}}{\simeq} A/B^{\otimes2}\right]
\]
compte tenu que l'on a un iso can
\[
\det(A)\simeq k
\]
et
\[
\deleted{\unsure{\det A\simeq}}\det B\otimes A/B\simeq\det A\ (\simeq k)
\]
Cet iso dépend d'une convention de signe,
% archive sheet number: 106
\scanpage{108}
\uline{Dém} L'équivalence de (i) et (ii) est valable, chaque fois qu'on a un fibré vectoriel orthogonal $A$ (rg \unsure{quelc.}) avec forme $d$ non dégénéré, et qu'on regarde \deleted{les} \inserted{un} \deleted{hyp} sous fibré vectoriel de codim $1$ (en supposant $\operatorname{rg}A$ pair \uline{ou} $2$ inv. dans $k$) -- \deleted{quand on \unsure{est}} \inserted{dans} ce dernier cas, inutile de parler de disc. \uline{divisé}, le discriminant ordinaire suffit).

En vertu de $\operatorname{disc}_{\varphi_B}=\operatorname{disc}_{\psi_B}=\inserted{2\operatorname{disc}'_{\det_{\mid B}}}$, \deleted{\unsure{la conv} et} la condition (ii\unsure{)} implique (ii$'$), et la réciproque est vraie si $2$ inv.

\deleted{Comme} Si $D$ est isotrope, il est orthogonal à lui-$=$, donc $D\unsure{\in|\subset}D^\perp=B$. L'inverse est vrai si $2$ est inv. (NB si $2=0$ dans $k$, alors on a \uline{tjrs} $D\subset B$ !); En effet, $D\subset B$ signifie que $D$ est orthogonal à lui-$=$, et donne $2d(x)=\varphi_d(x,x)=0$ Cela implique $2d(x)=0$ ou $d(x)=0$ pour $x\in D$ i.e. $D$ est isotrope.

On sait que les sous-alg. de Borel sont en corr. $1$-$1$ avec les $D$ isotropes \inserted{dans $A^{\unsure{s|\leftrightarrow|\natural}}$}, à \deleted{$D$ cor} une telle $D$
% archive sheet number: 107
\scanpage{109}
Les sous fibrés vectoriels $B$ de $A$, \inserted{contt $1$,} de rang $3$, correspondent par orthogonalité (pour $\varphi$ ou $\psi$, c'est kif kif) aux sous fibrés vectoriels de rang $1$ $D$ de $A^{\unsure{s|\leftrightarrow|\natural}}$.

D'autre part, pour un tel $B$, considérons
\[
\det_B=\det\mid B\ ,\quad\text{forme quadratique sur } B.
\]
elle admet un discriminant divisé
\[
\deleted{\unsure{\delta}}\operatorname{disc}'(\det_B)\quad\text{tels que}\quad2\operatorname{disc}'(\det_B)=\operatorname{disc}\psi_B
\]
où $\psi_B$ est la restriction de $\psi$ à $B$.

Ainsi
\[
\deleted{\unsure{2}}\operatorname{disc}\psi_B=\operatorname{disc}\varphi_B
=2\underbrace{\operatorname{disc}'(\det_B)}_{\substack{\parallel\,\text{déf}\\\delta(B)}}\unsure{=}
\]

\uline{Proposition} Les \inserted{quatre} conditions suivantes sont équivalentes

(i) $D$ est \uline{isotrope} i.e. $\det\mid D=0$

(ii) $\delta(B)=0$\qquad (\uline{NB} $\delta(B)\in\deleted{\unsure{(\det B)^{\otimes2}}}\det(B)^{\otimes(-2)}$).

\deleted{\unsure{(iii)}} (iii) $B$ est une sous-algèbre de $\unsure{A|M}$

% slanted
\margin{\deleted{Si $k$ un corps \unsure{sinon|il suffit} on sait pas si ça implique (iii$'$) \unsure{et}} $\longrightarrow$}

(iii$'$) \deleted{\unsure{(iii$'$)}} $B$ est une ``sous-alg. de Borel'' de $A$

Ces \deleted{\unsure{(iv)}} conditions impliquent les deux suivantes, et si $2.1_k$ est inv, \deleted{elles} ces dernières sont équivalentes aux précédentes

(iv) $D\subset B$

(\unsure{iv$'$|iii$'$|ii$'$}) \deleted{\unsure{$\operatorname{disc}=0$|$\operatorname{disc}$}} $\operatorname{disc}(\varphi_B)=0$\qquad ($\Longleftrightarrow\operatorname{disc}\psi_B=0$)
% archive sheet number: 108
\scanpage{110}
\uline{Cor.} Conditions équivalentes

(i) $D$ est anisotrope i.e. $\det_A\mid D$ est non nulle sur chaque fibre

(ii) $\delta(B)(=\delta(D))$ est un élémt inv. de $\deleted{\unsure{B^{\otimes-2}}}\inserted{(\det B)^{\otimes-2}}$

\deleted{i.e. $\det_{\unsure{A}}\mid D$ est partout non nulle}

i.e. $\det_B=\det_A\mid B$ est non dégénérée

(iii) $B$ n'est une sous-alg. sur aucune fibre

(iv) $B$ engendre $A$ comme \uline{algèbre} (pas comme module bien sûr)

\margin{\unsure{si} $2$ inv. sur $S$}
\[
\left\{\begin{aligned}
&(\mathrm{\unsure{v|iv}})\quad D\not\subset B\quad\text{sur tte fibre}\\
&(\mathrm{vi})\quad\varphi\mid B\text{ est non dégénérée}
\end{aligned}\right.
\]

Enfin, si $u,v$ sont deux éléments de $B$ tels que $1,u,v$ font une base, de $B$, ces conditions équivalent aussi à

(v) $1,u,v,uv$ forment une base de $A$

ou encore

(v$'$) $uv\bmod B$ est une \deleted{\unsure{section inv.}} \inserted{élémt inv.} de $A/B$

\begin{center}\rule{1.5cm}{0.35pt}\end{center}

\deleted{Supposons}

Nous regardons maintenant deux sous-fibrés \deleted{\unsure{vect.}} \deleted{\unsure{A}} $H,H'$ de rang $2$, contnt $1$ (ce sont néc. des sous-algèbres, par H.C.) et on suppose
\[
H\cap H'=k\qquad\text{sur tte fibre}
\]
donc
\[
H\oplus H'\overset{\text{déf.}}{=}B\qquad\text{sous-fibré de rang }3
\]
% archive sheet number: 109
\scanpage{111}
\deleted{\unsure{$H$},} Un $H\inserted{\supset1}$ est défini aussi par son orthogonal, et de $=$ pour $H'$
\[
V\subset A^{\unsure{s|\leftrightarrow|\natural}},\quad V'\subset A^s
\]
sous fibrés vectoriels de rang $2$, qui satisfont donc
\[
V+V'=A^s
\]
donc
\[
V\cap V'\overset{\text{def}}{=}D\qquad\unsure{\text{où}}
\]
est un sous fibré vect. de rang $1$, $\unsure{*}$

\deleted{Dire $H$ est étale} \deleted{que $H$ est}

Les $H$, correspondent à des sections \inserted{sur $k$ ou $k$-points} du \deleted{fibré} \inserted{plan} projectif \inserted{\unsure{duel|dual}} \deleted{\unsure{$V^\perp$}} $\check{\mathbb{P}}(A/k)$, et les $V$ à des droites dans le plan projectif dual, dans lequel il y a une conique canonique, celle des droites isotropes de $A^s$. Alors $H$ est \unsure{étale|étal|ét.} sur $k$ (\unsure{en tt que|les $H$ qui} algèbre sur $k$) sss en chaque \deleted{point} fibre \inserted{géom.}, la droite correspondante (polaire) $D_H=D_V$ est tangente à la conique, i.e. y découpe un diviseur \deleted{de} non réduit : un pt. Si cette condition est vérifiée pour
% slanted
\margin{\deleted{\uline{NB} Les $H$ en question sont en corr. $1$-$1$ avec les diviseurs de degré $2$ sur la conique qui sont non dégénérés.}}
% archive sheet number: 110
\scanpage{112}
$H$ et $H'$, on dit qu'ils sont \uline{en position générale} si sur chaque fibre géométrique, les deux paires de points sur la conique isotrope sont disjoints. Mais cela signifie \deleted{justement} qu'il n'y a pas de droite isotrope (sur la dite fibre) orth. \uline{à la fois} à $H$ et $H'$, i.e à $B$. Donc par notre proposition, cela signifie que $\delta_B$ est inv. \deleted{\unsure{i.e.|dire}.} que $\det\big|D$ est inversible. \deleted{\unsure{Si|Le}} Dans ce cas là, \deleted{\unsure{on choisit|choisit}} pour tt choix (possible localmt) \deleted{\unsure{Si on suppose $H,H'$ toutes deux|Si on suppose $H,H'$ étales|Supposons $H,H'$ étales}} de droites iso\unsure{~disjointes|. distinctes}
\[
 D_1,D_2\subset \deleted{\unsure{H}}V\cap\underset{\substack{\shortparallel\\\text{cône isotrope}}}{C},\qquad D'_1,D'_2\subset V'\cap C
\]
on a un rapport anharmonique
\[
 \lambda=\lambda_{(D_1,D_2)}(D'_1,D'_2)\in k
\]
\margin{$\lambda$ \unsure{partout} $\ne 0,1$ i.e $\lambda$ et $1-\lambda$ inv.}
qui est transformé en $1/\lambda$ si on echange $D_1$ et $D_2$ ou $D'_1$ et $D'_2$, et qui en dépend
% archive sheet number: 111
